\section{Online Monotone Metric Median}

In this section, we state and prove a general version of the
\(\texttt{sampled-lazy-medoid}\) algorithm. The result is not specific to
consensus clustering. It applies to any online median problem in a metric space
where the distances between the input trajectories are monotone over time.

Informally, we are given \(m\) input trajectories in a sequence of metric
spaces. At every time step \(t\), the algorithm must output a point whose total
distance to the \(m\) current input points is close to the optimal median cost.
The algorithm samples a logarithmic number of input trajectories and lazily
switches among them. With high probability, the sampled set contains a good
input point at every time step, and the lazy switching rule ensures low
recourse.

We first define the general problem.

\paragraph{Online monotone metric median.}
An instance consists of a sequence of metric spaces
\[
    (\mathcal M_t,d_t), \qquad t=1,\ldots,T,
\]
and \(m\) input trajectories
\[
    P_1,\ldots,P_m.
\]
At time \(t\), trajectory \(P_i\) is represented by a point
\[
    p_i^t\in \mathcal M_t.
\]
The algorithm must maintain an output point
\[
    S_t\in \mathcal M_t.
\]

For any point \(x\in\mathcal M_t\), define its objective value at time \(t\) as
\[
    \obj_t(x)
    :=
    \sum_{i=1}^m d_t(p_i^t,x).
\]
Let
\[
    \OPT_t
    \in
    \arg\min_{x\in\mathcal M_t}\obj_t(x)
\]
be an optimal median at time \(t\). Thus,
\[
    \obj_t(\OPT_t)
    =
    \min_{x\in\mathcal M_t}
    \sum_{i=1}^m d_t(p_i^t,x).
\]

We assume the instance is monotone in the following sense: for every pair of
input trajectories \(i,j\in[m]\), the distance
\[
    d_t(p_i^t,p_j^t)
\]
is nondecreasing in \(t\). That is, for all \(t<T\),
\[
    d_t(p_i^t,p_j^t)
    \le
    d_{t+1}(p_i^{t+1},p_j^{t+1}).
\]

For each input trajectory \(P_j\), define its medoid cost at time \(t\) by
\[
    c_j(t)
    :=
    \obj_t(p_j^t)
    =
    \sum_{i=1}^m d_t(p_i^t,p_j^t).
\]
We assume that medoid costs are bounded and discretized: for some parameter
\(\Lambda\),
\[
    c_j(t)\in \{0\}\cup [1,\Lambda]
\]
for every \(j\in[m]\) and every \(t\in[T]\).

The algorithm considered below always outputs one of the sampled input
trajectories. Therefore, in the abstract metric-median problem, we measure
recourse by the number of times the maintained sampled trajectory changes. If
the algorithm maintains index \(a_t\), then its switching recourse is
\[
    \left|
    \{t>1:a_t\neq a_{t-1}\}
    \right|.
\]

\begin{restatable}[Sampled lazy-medoid for online monotone metric median]{theorem}{metricmediansampledlazy}
\label{thm:metric-median-sampled-lazy}
Let \(\varepsilon,\eta\in(0,1]\), and let \(\delta\in(0,1)\). There is a randomized online algorithm for the online monotone metric median problem such that, with probability at least \(1-\delta\), for every time step \(t\in[T]\),
\[
    \obj_t(S_t)
    \le
    (2+\varepsilon)(1+\eta)\obj_t(\OPT_t).
\]
Moreover, on the same event, the switching recourse is at most
\[
    O\left(
        \frac{1}{\varepsilon\eta}
        \log\frac{T}{\delta}
        \log(1+\Lambda)
    \right).
\]
\end{restatable}

\subsection{Algorithm \texttt{sampled-lazy-medoid}}

In this section, we describe the algorithm
\(\texttt{sampled-lazy-medoid}\). The algorithm samples a small set of input
trajectories before the online process starts. It then maintains one sampled
trajectory and switches only when the current sampled trajectory becomes worse
than the best sampled trajectory by a multiplicative factor.

\paragraph{Input.}
The algorithm receives an online monotone metric median instance with \(m\)
input trajectories
\[
    P_1,\ldots,P_m
\]
and time horizon \(T\). It is also given parameters
\[
    \varepsilon,\eta\in(0,1],
    \qquad
    \delta\in(0,1).
\]

\paragraph{Output.}
At every time step \(t\), the algorithm outputs a point
\[
    S_t\in\mathcal M_t.
\]

\paragraph{Sampling step.}
Before the online process starts, the algorithm samples
\[
    q
    =
    \left\lceil
        \frac{2+\varepsilon}{\varepsilon}
        \ln\frac{T}{\delta}
    \right\rceil
\]
indices independently and uniformly from \([m]\). Let
\[
    K
\]
be the resulting sampled multiset of indices.

For every sampled index \(j\in K\), define
\[
    c_j(t)
    =
    \obj_t(p_j^t)
    =
    \sum_{i=1}^m d_t(p_i^t,p_j^t).
\]
Also define the best sampled medoid cost
\[
    M_K(t)
    =
    \min_{j\in K} c_j(t).
\]

\paragraph{Lazy switching rule.}
At time \(t=1\), the algorithm chooses
\[
    a_1\in \arg\min_{j\in K} c_j(1),
\]
breaking ties lexicographically, and outputs
\[
    S_1=p_{a_1}^1.
\]

At every later time step \(t>1\), after the points
\[
    p_1^t,\ldots,p_m^t
\]
are revealed, the algorithm computes \(c_j(t)\) for all \(j\in K\).

If the previously maintained trajectory \(a_{t-1}\) satisfies
\[
    c_{a_{t-1}}(t)
    \le
    (1+\eta)M_K(t),
\]
then the algorithm keeps the same trajectory:
\[
    a_t=a_{t-1}.
\]
Otherwise, it switches to a sampled minimizer:
\[
    a_t\in \arg\min_{j\in K} c_j(t),
\]
again breaking ties lexicographically.

Finally, the algorithm outputs
\[
    S_t=p_{a_t}^t.
\]

\subsection{Competitive Ratio and Recourse of \texttt{sampled-lazy-medoid}}

In this section, we prove \cref{thm:metric-median-sampled-lazy}. The proof has two parts. First, we show that with high probability, the sampled set \(K\) contains a near-best input trajectory at every time step. Second, we use monotonicity to prove that the lazy switching rule changes the maintained trajectory only polylogarithmically many times.

\begin{lemma}[Average medoid bound]
\label{lem:metric-average-medoid}
For every time step \(t\),
\[
    \frac1m
    \sum_{j=1}^m c_j(t)
    \le
    2\obj_t(\OPT_t).
\]
\end{lemma}

\begin{proof}
By definition,
\[
\begin{aligned}
    \frac1m\sum_{j=1}^m c_j(t)
    &=
    \frac1m
    \sum_{j=1}^m
    \sum_{i=1}^m
    d_t(p_i^t,p_j^t).
\end{aligned}
\]
For every pair \(i,j\), the triangle inequality gives
\[
    d_t(p_i^t,p_j^t)
    \le
    d_t(p_i^t,\OPT_t)+d_t(\OPT_t,p_j^t).
\]
Therefore,
\[
\begin{aligned}
    \frac1m
    \sum_{j=1}^m
    \sum_{i=1}^m
    d_t(p_i^t,p_j^t)
    &\le
    \frac1m
    \sum_{j=1}^m
    \sum_{i=1}^m
    \left(
        d_t(p_i^t,\OPT_t)+d_t(\OPT_t,p_j^t)
    \right)\\
    &=
    \frac1m
    \left(
        m\sum_{i=1}^m d_t(p_i^t,\OPT_t)
        +
        m\sum_{j=1}^m d_t(\OPT_t,p_j^t)
    \right)\\
    &=
    2\sum_{i=1}^m d_t(p_i^t,\OPT_t)\\
    &=
    2\obj_t(\OPT_t).
\end{aligned}
\]
This proves the lemma.
\end{proof}

\begin{lemma}[The sample contains a good medoid]
\label{lem:metric-sample-good}
With probability at least \(1-\delta\), for every time step \(t\in[T]\),
\[
    M_K(t)
    \le
    (2+\varepsilon)\obj_t(\OPT_t).
\]
\end{lemma}

\begin{proof}
Fix a time step \(t\). Define the set of good indices
\[
    G_t
    =
    \left\{
        j\in[m]:
        c_j(t)
        \le
        (2+\varepsilon)\obj_t(\OPT_t)
    \right\}.
\]

We first show that \(G_t\) contains a nontrivial fraction of all indices. By
\cref{lem:metric-average-medoid},
\[
    \frac1m\sum_{j=1}^m c_j(t)
    \le
    2\obj_t(\OPT_t).
\]
If fewer than
\[
    \frac{\varepsilon}{2+\varepsilon}m
\]
indices were in \(G_t\), then more than
\[
    \frac{2}{2+\varepsilon}m
\]
indices would have cost strictly larger than
\[
    (2+\varepsilon)\obj_t(\OPT_t).
\]
This would imply
\[
    \frac1m\sum_{j=1}^m c_j(t)
    >
    \frac{2}{2+\varepsilon}
    \cdot
    (2+\varepsilon)\obj_t(\OPT_t)
    =
    2\obj_t(\OPT_t),
\]
contradicting the average bound. Hence
\[
    |G_t|
    \ge
    \frac{\varepsilon}{2+\varepsilon}m.
\]

Each sampled index is chosen uniformly from \([m]\). Therefore, the probability that none of the \(q\) samples lies in \(G_t\) is at most
\[
    \left(
        1-\frac{\varepsilon}{2+\varepsilon}
    \right)^q
    \le
    \exp\left(
        -\frac{\varepsilon}{2+\varepsilon}q
    \right).
\]
By the choice of
\[
    q
    =
    \left\lceil
        \frac{2+\varepsilon}{\varepsilon}
        \ln\frac{T}{\delta}
    \right\rceil,
\]
this failure probability is at most
\[
    \frac{\delta}{T}.
\]

Taking a union bound over all time steps \(t=1,\ldots,T\), with probability at least \(1-\delta\), the sampled set \(K\) intersects \(G_t\) for every \(t\). On this event, for every time \(t\),
\[
    M_K(t)
    =
    \min_{j\in K} c_j(t)
    \le
    (2+\varepsilon)\obj_t(\OPT_t).
\]
This proves the lemma.
\end{proof}

\begin{lemma}[Competitive ratio]
\label{lem:metric-sampled-competitive}
On the event of \cref{lem:metric-sample-good}, for every time step \(t\),
\[
    \obj_t(S_t)
    \le
    (2+\varepsilon)(1+\eta)\obj_t(\OPT_t).
\]
\end{lemma}

\begin{proof}
By construction, the maintained index \(a_t\) always satisfies
\[
    c_{a_t}(t)
    \le
    (1+\eta)M_K(t).
\]
Indeed, if the previous index \(a_{t-1}\) satisfies this inequality, the algorithm keeps it. Otherwise, the algorithm switches to an index attaining \(M_K(t)\), in which case the inequality holds trivially.

Since
\[
    S_t=p_{a_t}^t,
\]
we have
\[
    \obj_t(S_t)=c_{a_t}(t).
\]
Therefore,
\[
    \obj_t(S_t)
    \le
    (1+\eta)M_K(t).
\]
On the event of \cref{lem:metric-sample-good},
\[
    M_K(t)
    \le
    (2+\varepsilon)\obj_t(\OPT_t).
\]
Combining the two inequalities,
\[
    \obj_t(S_t)
    \le
    (2+\varepsilon)(1+\eta)\obj_t(\OPT_t).
\]
\end{proof}

We now prove the recourse bound.

\begin{lemma}[Monotonicity of sampled medoid costs]
\label{lem:metric-costs-monotone}
For every sampled index \(j\in K\), the function \(c_j(t)\) is nondecreasing in \(t\). Moreover, \(M_K(t)\) is nondecreasing in \(t\).
\end{lemma}

\begin{proof}
Fix \(j\in K\). By definition,
\[
    c_j(t)
    =
    \sum_{i=1}^m d_t(p_i^t,p_j^t).
\]
By the monotonicity assumption, for every \(i\in[m]\),
\[
    d_t(p_i^t,p_j^t)
    \le
    d_{t+1}(p_i^{t+1},p_j^{t+1}).
\]
Summing over \(i\), we get
\[
    c_j(t)\le c_j(t+1).
\]
Thus \(c_j(t)\) is nondecreasing.

Since \(M_K(t)\) is the minimum of the nondecreasing functions \(c_j(t)\), \(j\in K\), it is also nondecreasing.
\end{proof}

\begin{lemma}[Switching bound]
\label{lem:metric-switching-bound}
The total number of switches made by \(\texttt{sampled-lazy-medoid}\) is at most
\[
    O\left(
        |K|\cdot
        \frac{1}{\eta}
        \log(1+\Lambda)
    \right).
\]
\end{lemma}

\begin{proof}
Define the smoothed sampled minimum
\[
    \widehat M_K(t)
    :=
    1+M_K(t).
\]
Since every medoid cost lies in \(\{0\}\cup[1,\Lambda]\), we have
\[
    1\le \widehat M_K(t)\le 1+\Lambda
\]
for every \(t\).

Fix a sampled index \(j\in K\). We bound how many times the algorithm can switch back to \(j\) after previously switching away from \(j\).

Suppose the algorithm switches away from \(j\) at time \(t\). Then, by the switching rule,
\[
    c_j(t)>(1+\eta)M_K(t).
\]
Suppose that at some later time \(t'>t\), the algorithm switches back to \(j\). Since the algorithm switches to a sampled minimizer,
\[
    c_j(t')=M_K(t').
\]
By monotonicity of \(c_j\),
\[
    c_j(t')\ge c_j(t).
\]
Hence
\[
    M_K(t')
    =
    c_j(t')
    \ge
    c_j(t)
    >
    (1+\eta)M_K(t).
\]

We now show that the smoothed quantity \(\widehat M_K\) grows by a multiplicative factor. If \(M_K(t)=0\), then
\[
    c_j(t)>0.
\]
Since medoid costs are either zero or at least \(1\), this implies
\[
    c_j(t)\ge 1.
\]
Therefore,
\[
    M_K(t')=c_j(t')\ge 1,
\]
and hence
\[
    \widehat M_K(t')=1+M_K(t')\ge 2.
\]
Since \(\widehat M_K(t)=1\), we get
\[
    \widehat M_K(t')
    \ge
    2\widehat M_K(t)
    \ge
    \left(1+\frac{\eta}{2}\right)\widehat M_K(t),
\]
because \(\eta\le 1\).

Now suppose \(M_K(t)\ge 1\). Then
\[
\begin{aligned}
    \widehat M_K(t')
    &=1+M_K(t')\\
    &>
    1+(1+\eta)M_K(t).
\end{aligned}
\]
Since \(M_K(t)\ge 1\),
\[
    1+(1+\eta)M_K(t)
    \ge
    \left(1+\frac{\eta}{2}\right)(1+M_K(t)).
\]
Thus,
\[
    \widehat M_K(t')
    \ge
    \left(1+\frac{\eta}{2}\right)\widehat M_K(t).
\]

Therefore, between two consecutive selections of the same sampled index \(j\), after the algorithm has switched away from \(j\), the smoothed sampled minimum \(\widehat M_K\) increases by a factor at least
\[
    1+\frac{\eta}{2}.
\]

Since
\[
    \widehat M_K(t)\le 1+\Lambda
\]
at all times, each sampled index can be reselected at most
\[
    O\left(
        \log_{1+\eta/2}(1+\Lambda)
    \right)
    =
    O\left(
        \frac{1}{\eta}
        \log(1+\Lambda)
    \right)
\]
times. Summing over all sampled indices \(j\in K\), the total number of switches is at most
\[
    O\left(
        |K|\cdot
        \frac{1}{\eta}
        \log(1+\Lambda)
    \right).
\]
\end{proof}

\begin{lemma}[Recourse bound]
\label{lem:metric-recourse-bound}
The switching recourse of \(\texttt{sampled-lazy-medoid}\) is at most
\[
    O\left(
        \frac{1}{\varepsilon\eta}
        \log\frac{T}{\delta}
        \log(1+\Lambda)
    \right).
\]
\end{lemma}

\begin{proof}
By definition, the switching recourse is the number of times the maintained sampled index changes.

By \cref{lem:metric-switching-bound}, the total number of switches is at most
\[
    O\left(
        |K|\cdot
        \frac{1}{\eta}
        \log(1+\Lambda)
    \right).
\]
The sample size is
\[
    |K|
    =
    q
    =
    O\left(
        \frac{1}{\varepsilon}
        \log\frac{T}{\delta}
    \right).
\]
Therefore, the switching recourse is at most
\[
    O\left(
        \frac{1}{\varepsilon\eta}
        \log\frac{T}{\delta}
        \log(1+\Lambda)
    \right).
\]
\end{proof}

\begin{proof}[Proof of \cref{thm:metric-median-sampled-lazy}]
By \cref{lem:metric-sample-good}, with probability at least \(1-\delta\), for every time step \(t\),
\[
    M_K(t)
    \le
    (2+\varepsilon)\obj_t(\OPT_t).
\]
Condition on this event.

By \cref{lem:metric-sampled-competitive}, for every time step \(t\),
\[
    \obj_t(S_t)
    \le
    (2+\varepsilon)(1+\eta)\obj_t(\OPT_t).
\]
By \cref{lem:metric-recourse-bound}, the switching recourse is at most
\[
    O\left(
        \frac{1}{\varepsilon\eta}
        \log\frac{T}{\delta}
        \log(1+\Lambda)
    \right).
\]
This proves the theorem.
\end{proof}

\begin{restatable}[Consensus clustering instantiation]{corollary}{sampledlazymedoidconsensus}
\label{cor:sampled-lazy-consensus}
For online consensus clustering, there is a randomized online algorithm that, with probability at least \(1-\delta\), satisfies for every time step \(t\),
\[
    \obj_t(S_t)
    \le
    (2+\varepsilon)(1+\eta)\obj_t(\OPT_t),
\]
and every vertex has recourse at most
\[
    O\left(
        \frac{1}{\varepsilon\eta}
        \log\frac{n}{\delta}
        \log(mn)
    \right).
\]
In particular, by setting
\[
    \varepsilon=\eta=\frac{1}{\log(mn)}
\]
and
\[
    \delta=(mn)^{-\Theta(1)},
\]
we obtain a \((2+o(1))\)-competitive algorithm with
\[
    O(\log^4(mn))
\]
recourse, with high probability.
\end{restatable}

\begin{proof}
We instantiate the online monotone metric median theorem with
\[
    \mathcal M_t=\{\text{clusterings of }V^t\},
\]
and
\[
    d_t(\m A,\m B)=\dist_{V^t}(\m A,\m B).
\]
The \(m\) input trajectories are
\[
    p_i^t=\m C_i^t,
\]
the restrictions of the \(m\) input clusterings to the arrived vertex set.

The objective becomes
\[
    \obj_t(\m C)
    =
    \sum_{i=1}^m
    \dist_{V^t}(\m C_i^t,\m C),
\]
which is exactly the online consensus clustering objective.

The monotonicity condition holds because when a new vertex arrives, the distance between any two restricted input clusterings only gains additional nonnegative pairwise disagreement terms. Thus, for every \(i,j\),
\[
    \dist_{V^t}(\m C_i^t,\m C_j^t)
    \le
    \dist_{V^{t+1}}(\m C_i^{t+1},\m C_j^{t+1}).
\]

Moreover, all consensus objectives are integer-valued and bounded by
\[
    \Lambda
    =
    m\binom n2.
\]
The time horizon is
\[
    T=n.
\]
Applying \cref{thm:metric-median-sampled-lazy}, we get that, with probability at least \(1-\delta\), for every time step \(t\),
\[
    \obj_t(S_t)
    \le
    (2+\varepsilon)(1+\eta)\obj_t(\OPT_t),
\]
and the switching recourse is at most
\[
    O\left(
        \frac{1}{\varepsilon\eta}
        \log\frac{n}{\delta}
        \log(mn)
    \right).
\]

Finally, in online consensus clustering, while the algorithm follows the same sampled input clustering, no previously arrived vertex changes cluster ID. A vertex can change cluster ID only when the algorithm switches from one sampled input clustering to another. Therefore the vertex recourse is at most the switching recourse.

Setting
\[
    \varepsilon=\eta=\frac{1}{\log(mn)}
\]
and
\[
    \delta=(mn)^{-\Theta(1)}
\]
gives
\[
    (2+\varepsilon)(1+\eta)
    =
    2+o(1),
\]
and
\[
    O\left(
        \frac{1}{\varepsilon\eta}
        \log\frac{n}{\delta}
        \log(mn)
    \right)
    =
    O(\log^4(mn)).
\]
This proves the corollary.
\end{proof}